\section{Coding 是 AGI 第二幕}

访谈中最强的判断，是 Coding 把 AI 从 Chatbot 第一幕推向 Agent 第二幕。Chatbot 的核心是问答、搜索、总结和对话；Coding 的核心是把自然语言意图转成可执行方案，让模型直接改代码、跑任务、调工具、处理数据和构建系统。它让 AI 从“描述世界”进入“改造数字世界”。

\begin{figure}[H]
\centering
\includegraphics[width=0.94\textwidth]{coding-second-act.png}
\caption{Coding 是 AGI 第二幕：从聊天到干活，再到模型 OS。}
\end{figure}

\begin{knowledgebox}{读图：为什么 Coding 是转折点}
图中第一幕是 Chatbot，价值来自信息压缩；第二幕是 Coding/Agent，价值来自完成工作；第三幕是 Model as OS，模型成为用户意图到应用执行之间的调度层。Coding 的关键不是写代码本身，而是代码能表达数字世界的大多数解决方案。
\end{knowledgebox}

\subsection{“语言即世界，代码即方案”}

访谈里有一句很适合作为本章标题的话：自然语言是对世界的描述，代码是对 solution 的描述。自然语言能表达意图，代码能表达执行路径。若模型能稳定写代码，它不仅能回答问题，还能操作文件、调用 API、处理数据、自动化流程、构建产品。这也是为什么 Coding 模型被类比为 GPU：它不是单个应用，而是加速器。

\begin{importantbox}{Coding 作为加速器}
GPU 加速模型训练，Coding 模型加速人类和 AI 的研究/工程循环。一个想法从两三周跑通变成两三天，意味着实验吞吐、产品迭代和研究反馈都被压缩。
\end{importantbox}

\subsection{研究员不再亲自写代码}

广密提到，前沿实验室研究员和强程序员已经大量减少亲自写代码，转向“AI 写，人来审”。如果 AI 能达到 CTO 或首席架构师级别的实现能力，人类工作会从写实现转向定方向、审设计、验结果和承担责任。这和 EP139 的 Agent 生产化、EP138 的后训练系统形成呼应。

\begin{warningbox}{不要把“人不写代码”理解成“人不重要”}
人类不写大量实现代码，不等于退出工程。真正稀缺的能力会转向定义问题、识别错误、设计实验、组织系统和判断价值。审查能力、架构能力和任务定义能力反而更重要。
\end{warningbox}

\subsection{本章小结}

Coding 之所以是第二幕，是因为它把 AI 从语言输出变成行动系统。它不仅提升程序员效率，也会加速 AI 研究本身，成为 AGI 路线中的核心放大器。

